\documentclass[12pt]{article}
\usepackage{amsmath,amssymb}

\begin{document}
\section*{{2023 AIME I Problems/Problem 7}}
Source: AoPS Wiki

\subsection*{Solution}
$n$ can either be $0$ or $1$ (mod $2$ ). Case 1: $n \equiv 0 \pmod{2}$ Then, $n \equiv 2 \pmod{4}$ , which implies $n \equiv 1 \pmod{3}$ and $n \equiv 4 \pmod{6}$ , and therefore $n \equiv 3 \pmod{5}$ . Using CRT , we obtain $n \equiv 58 \pmod{60}$ , which gives $16$ values for $n$ . Case 2: $n \equiv 1 \pmod{2}$ $n$ is then $3 \pmod{4}$ . If $n \equiv 0 \pmod{3}$ , $n \equiv 3 \pmod{6}$ , a contradiction. Thus, $n \equiv 2 \pmod{3}$ , which implies $n \equiv 5 \pmod{6}$ . $n$ can either be $0 \pmod{5}$ , which implies that $n \equiv 35 \pmod{60}$ by CRT, giving $17$ cases; or $4 \pmod{5}$ , which implies that $n \equiv 59 \pmod{60}$ by CRT, giving $16$ cases. The total number of extra-distinct numbers is thus $16 + 16 + 17 = \boxed{049}$ . ~mathboy100

\end{document}
